\chapter{Lecture 18: Antiderivatives of more complicated functions I}
In section \ref{derivativetable1} we developed a small table of antiderivative formulas. Then, in section \ref{Using FTC to find definite integrals}, we learned how to use antiderivatives to find definite integrals.  \\

\noindent In section \ref{product rule} we learned a formula (the ``product rule'') for differentiating functions that can be expressed as a product of two functions. Then, in section \ref{chain rule}, we learned a formula (the chain rule) for differentiating functions which can be expressed as a composition of two functions. These were
\[
\frac{d}{dx}\brac{f(x)g(x)} = f'(x)g(x) + f(x)g'(x) \qquad \text{and} \qquad \frac{d}{dx}\brac{f(g(x))} = f'(g(x))g'(x)
\]

\noindent There are two corresponding antiderivative formulas, one for antidifferentiating products of functions (called \textbf{integration by parts}) and one for antidifferentiating compositions of functions (called \textbf{integration by substitution}); these can be derived from the product and chain rules respectively. \\

\noindent In this lecture we will learn to use integration by parts, and in the next lecture we will learn integration by substitution.

\section{Integration by parts} \label{integration by parts}
Starting from the product rule for derivatives, we can derive the integration by parts rule. The product rule is
\[
\frac{d}{dx}\brac{f(x)g(x)} = f'(x)g(x) + f(x)g'(x)
\]
%Define a new function, $f(x)$, so that $f(x) = k'(x)$. Then $k(x) = \int f(x) \; dx = F(x)$. So we can rewrite the above as
%\[
%\frac{d}{dx}\brac{F(x)g(x)} = f(x)g(x) + F(x)g'(x)
%\]
%Antidifferentiating both sides gives
%\[
%\int\frac{d}{dx}\brac{F(x)g(x)} \; dx = \int f(x)g(x) + F(x)g'(x) \; dx
%\]
%The antiderivative cancels the derivative on the left hand side, and on the right hand side we use linearity, giving
%\[
%F(x)g(x) = \int f(x)g(x) \; dx + \int F(x)g'(x) \; dx
%\]
%Subtracting $ \int F(x)g'(x) \; dx$ from both sides gives us the integration by parts rule:
%
%  	\bigskip
%	
%	\hfil\fbox{
%	\begin{minipage}{0.53\columnwidth}{
%	
%	\vskip 0.02cm
%	\begin{center}
%	$\displaystyle{ \int f(x)g(x) \; dx = F(x)g(x) - \int F(x)g'(x) \; dx}$
%	\end{center}
%	}
%	\end{minipage}}
%	\bigskip
Antidifferentiating both sides gives
\[
\int\frac{d}{dx}\brac{f(x)g(x)} \; dx = \int f'(x)g(x) + f(x)g'(x) \; dx
\]
The antiderivative cancels the derivative on the left hand side, and on the right hand side we use linearity, giving
\[
f(x)g(x) = \int f'(x)g(x) \; dx + \int f(x)g'(x) \; dx
\]
Subtracting $ \int f(x)g'(x) \; dx$ from both sides gives us the integration by parts rule:

  	\bigskip
	
	\hfil\fbox{
	\begin{minipage}{0.53\columnwidth}{
	
	\vskip 0.02cm
	\begin{center}
	$\displaystyle{ \int f'(x)g(x) \; dx = f(x)g(x) - \int f(x)g'(x) \; dx}$
	\end{center}
	}
	\end{minipage}}
	\bigskip

\noindent Another way that it is commonly written is to let $u = f(x)$, and $v = g(x)$; then

  	\bigskip
	
	\hfil\fbox{
	\begin{minipage}{0.4\columnwidth}{
	
	\vskip 0.02cm
	\begin{center}
	$\displaystyle{ \int u'v \; dx = uv - \int uv' \; dx}$
	\end{center}
	}
	\end{minipage}}
	\bigskip

\noindent When using the integration by parts rule, we need to be unusually mindful. If we want to antidifferentiate a product, then we need to nominate one factor to be $f'(x)$, and another factor to be $g(x)$. Looking at the right hand side, this means that we will need to integrate $f'(x)$ to get $f(x)$. Also, we will need to differentiate $g(x)$ to get $g'(x)$. We then need the expression $\int f(x)g'(x) \; dx$ on the far right to be something simpler that can be integrated easier than the expression that we started with; otherwise we end up having to use integration by parts again on $\int f(x)g'(x) \; dx$, which can lead to us being stuck in an endless loop for all eternity. \\

\noindent It is perhaps best taught by example. \\

\noindent A good function to practice on is $h(x) = x\e^{-x}$. Suppose that we haphazardly select $f'(x) = x$ and $g(x) = \e^{-x}$. Then
\[
f(x) = \int x \; dx = \frac{1}{2}x^2 \qquad \text{and} \qquad g'(x) = -\e^{-x}
\]
Now, using integration by parts
\begin{align*}
\int x\e^{-x} \; dx & = f(x)g(x) - \int f(x)g'(x) \; dx \\
		        & = \frac{1}{2}x^2\e^{-x} - \int \frac{1}{2}x^2 (-\e^{-x}) \; dx. \tag{plugging in our results} \\
		        & = \frac{1}{2}x^2\e^{-x} + \frac{1}{2} \int x^2 \e^{-x} \; dx. \tag{linearity}
\end{align*}
We started off needing to integrate $x\e^{-x}$, and now we need to integrate $x^2\e^{-x}$. Things have not gotten any better; they have gotten worse. If we keep going like this, we may never make it out. \\

\noindent We made a mistake, though it is easily fixed by changing our approach. This time we let $f'(x) = \e^{-x}$, and we let $g(x) = x$. Then
\[
f(x) = \int \e^{-x} \; dx = -\e^{-x} \qquad \text{and} \qquad g'(x) = 1
\]
This time, using integration by parts gives
\begin{align*}
\int x\e^{-x} \; dx & = f(x)g(x) - \int f(x)g'(x) \; dx \\
		& = (-\e^{-x})x - \int (-\e^{-x})\times1 \; dx \tag{plugging in our results} \\
		& = -\e^{-x}x + \int \e^{-x} \; dx \tag{linearity} \\
		& = -\e^{-x}x - \e^{-x} \\
		& = -\e^{-x}(x+1).
\end{align*}
This time things went much more smoothly. The trick was to recognise that we would eventually need to integrate a product involving $g'(x)$; letting $g(x) = x$ meant that we would (conveniently) end up integrating a product involving $g'(x) = 1$. 

\subsection{Example}
Find an antiderivative of $x\ln(x)$. \\

\solution
Since we do not even know how to integrate $\ln(x)$ yet (but we do know how to differentiate it), we decide to choose $g(x) = \ln(x)$ and $f'(x) = x$. Then
\[
g'(x) = \frac{1}{x} \qquad \text{and} \qquad f(x) = \int x \; dx = \frac{1}{2}x^2 
\]
Now, using integration by parts
\begin{align*}
\int x\ln(x) \; dx & = f(x)g(x) - \int f(x)g'(x) \; dx	\\
		& = \frac{1}{2}x^2\ln(x) - \int \frac{1}{2}x^2 \frac{1}{x} \; dx \tag{plugging in our results} \\
		& =  \frac{1}{2}x^2\ln(x) - \int \frac{1}{2}x \; dx \\
		& = \frac{1}{2}x^2\ln(x) - \frac{1}{4}x^2 \\
		& = \frac{1}{2}x^2\brac{\ln(x) - \frac{1}{2}}.
\end{align*}

\subsection{Antiderivative of $\boldsymbol{\ln(x)}$}
So far, we know that $\frac{d}{dx}\brac{\ln(x)} = \frac{1}{x}$. However, our table of antiderivatives (given in section \ref{derivativetable1}) does not include a formula for the antiderivative of $\ln(x)$. We now develop one using a clever application of integration by parts. \\

\noindent If we let $g(x) = \ln(x)$, and $f'(x) = 1$, then $\ln(x) = f'(x)g(x)$; this means that we can use integration by parts on $\ln(x)$. We also need the following:
\[
f(x) = \int 1 \; dx = x \qquad \text{and} \qquad g'(x) = \frac{1}{x}.
\]
Using integration by parts we have
\begin{align*}
\int \ln(x) \; dx & = f(x)g(x) - \int f(x)g'(x) \; dx \\
	& = x\ln(x) - \int x\frac{1}{x} \; dx \\
	& = x\ln(x) - \int 1 \; dx \\
	& = x\ln(x) - x.
\end{align*}
We have found a useful antiderivative formula:

  	\bigskip
	
	\hfil\fbox{
	\begin{minipage}{0.35\columnwidth}{
	
	\vskip 0.02cm
	\begin{center}
	$\displaystyle{\int \ln(x) \; dx = x\ln(x) - x}$
	\end{center}
	}
	\end{minipage}}
	\bigskip

\section{Integration by parts with definite integrals} \label{sec18.2}
When using integration by parts on a definite integral, we need to be mindful to think about the terminals. The rule becomes

  	\bigskip
	
	\hfil\fbox{
	\begin{minipage}{0.63\columnwidth}{
	
	\vskip 0.02cm
	\begin{center}
	$\displaystyle{ \int_a^b f'(x)g(x) \; dx = \squac{f(x)g(x)}_a^b - \int_a^b f(x)g'(x) \; dx}$
	\end{center}
	}
	\end{minipage}}
	\bigskip

\noindent For example, suppose we wish to evaluate $\int_0^\infty x\e^{-x} \; dx$. In section \ref{integration by parts} we used integration by parts to find that
\[
\int x\e^{-x} \; dx =  -\e^{-x}x + \int \e^{-x} \; dx 
\]
Adding the terminals we get
\begin{align*}
\int_0^\infty x\e^{-x} \; dx & =  \squac{-\e^{-x}x}_0^\infty + \int_0^\infty \e^{-x} \; dx  \\
		& = 0 + \int_0^\infty \e^{-x} \; dx \tag{since $\e^{-x} \rightarrow 0$ as $x \rightarrow \infty$} \\
		& = \squac{-\e^{-x}}_0^\infty \\
		& = 0 - (-\e^{-0}) \tag{since $\e^{-x} \rightarrow 0$ as $x \rightarrow \infty$} \\
		& = 1. \tag{as $\e^0 = 1$ (index laws)}
\end{align*}

\subsection{Example}
Find $\int_0^3 x^2\e^{-x} \; dx$. \\

\solution
Let $f'(x) = \e^{-x}$ and let $g(x) = x^2$. Then
\[
f(x) = \int \e^{-x} \; dx = -\e^{-x} \qquad \text{and} \qquad g'(x) = 2x
\]
\noindent Using integration by parts,
\begin{align*}
\int_0^3 x^2\e^{-x} \; dx & = \squac{f(x)g(x)}_0^3 - \int_0^3 f(x)g'(x) \; dx \\
	& = \squac{-\e^{-x}x^2}_0^3 - \int_0^2 (-\e^{-x})2x \; dx \tag{plugging in our results} \\
	& = \squac{-\e^{-x}x^2}_0^3 + 2\int_0^2 \e^{-x}x \; dx \tag{linearity} \\
	& = \squac{-\e^{-x}x^2}_0^3 + 2\squac{-\e^{-x}(x+1)}_0^3 \tag{by the result in section \ref{integration by parts}} \\
	& = -\e^{-3}3^2 - 0 + 2\brac{-\e^{-3}(3+1) - (-1)} \\
	& = 2 - 17\e^{-3}.
\end{align*}

\section{Proofs for some properties of the gamma function} \label{sec18.1}
\subsubsection{Proof that $\boldsymbol{\Gamma(\alpha) = (\alpha-1)\Gamma(\alpha-1)}$} 
In section \ref{gamma properties} we mentioned this important property of the gamma function. Now we can prove it using integration by parts: \\

\noindent By definition, the gamma function is
\[
\Gamma(\alpha) = \int_0^\infty x^{\alpha - 1}\e^{-x} \; dx
\]
Let $f'(x) = \e^{-x}$ and let $g(x) = x^{\alpha - 1}$. Then
\[
f(x) = \int \e^{-x} \; dx = -\e^{-1} \qquad \text{and} \qquad g'(x) = (\alpha - 1)x^{\alpha - 2}.
\]
Now, using integration by parts gives
\begin{align*}
\Gamma(\alpha) & = \squac{f(x)g(x)}_0^\infty - \int_0^\infty f(x)g'(x) \; dx \\
	& = \squac{-\e^{-x}x^{\alpha - 1}}_0^\infty + (\alpha - 1)\int_0^\infty \e^{-x}x^{\alpha - 2} \; dx \tag{using linearity}\\
	& = 0 + (\alpha - 1)\int_0^\infty e^{-x}x^{\alpha - 2} \; dx \tag{since $\e^{-x} \rightarrow 0$ as $x \rightarrow \infty$} \\
	& = (\alpha - 1)\Gamma(\alpha - 1). \tag{definition of gamma function}
\end{align*}
Which is what we wanted to show.

\subsubsection{Proofs that $\boldsymbol{\Gamma(1) = 1}$ and $\boldsymbol{\Gamma(n) = (n-1)!}$ for $\boldsymbol{n \in \N}$}
By the definition of the gamma function we have
\[
\Gamma(2) = \int_0^\infty x^{2-1}\e^{-x} \; dx = \int_0^\infty x\e^{-x} \; dx
\]
We evaluated this integral in \ref{sec18.2}; we found that $\int_0^\infty x\e^{-x} \; dx = 1$. Hence,
\[
\Gamma(2) = 1.
\]
Using the fact that 
\[
\boldsymbol{\Gamma(\alpha) = (\alpha - 1)\Gamma(\alpha - 1)} \tag{$\bigstar$}
\]
we have
\[
\Gamma(2) = (2-1)\Gamma(2-1) = \Gamma(1)
\]
So $\Gamma(1) = 1$ as well. \\

\noindent Now, let $n$ be a natural number. For the following steps we just use $\bigstar$ repeatedly:
\begin{align*}
\Gamma(n) & = (n-1)\Gamma(n-1)  \tag{by $\bigstar$} \\
	& = (n-1)(n-2)\Gamma(n-2) \tag{as $\Gamma(n-1) = (n-2)\Gamma(n-2)$} \\
	& = (n-1)(n-2)(n-3)\Gamma(n-3)  \tag{as $\Gamma(n-2) = (n-3)\Gamma(n-3)$}\\
	& = \ldots \tag{continue applying $\bigstar$ repeatedly} \\
	& = (n-1)(n-2)(n-3)\times\ldots\times\Gamma(3) \\
	& = (n-1)(n-2)(n-3)\times\ldots\times2\Gamma(2) \tag{as $\Gamma(3) = 2\Gamma(2)$} \\
	& = (n-1)(n-2)(n-3)\times\ldots\times2\times1\times\Gamma(1). \tag{as $\Gamma(2) = 1\Gamma(1)$} \\
	& = (n-1)!  \tag{as $\Gamma(1) = 1$}
\end{align*}
Hence $\Gamma(n) = (n-1)!$, as required.

